\section{Checked Lean nucleus and a concrete formalization route}
\label{sec:formalization}

The archive separates an executable, kernel-checked development from the
larger object-language implementation that remains to be written. This
section records that boundary and gives the definitions and induction
targets needed to cross it. The detailed module plan is also supplied in
\path{lean/OBJECT_LANGUAGE_PLAN.md}.

\subsection{Exactly what the delivered Lean development proves}

\path{lean/QuantitativeKernel.lean} contains 19 theorem declarations with
complete proofs. It and \path{lean/Examples.lean} were checked using Lean
4.32.0, compiler commit
\texttt{8c9756b28d64dab099da31a4c09229a9e6a2ef35}.
There are no \texttt{sorry} placeholders and no new axiom declarations.
The complete replay output is in \path{lean/verification.log}.

\begin{center}
\small
\begin{tabular}{p{0.27\textwidth}p{0.63\textwidth}}
\toprule
Delivered component & Exact status and scope\\
\midrule
\path{positive_transfer} & Checked host Lean theorem: preservation of positive
formulas under an atomic transfer contract and explicit forward/backward
witness selectors. No axiom dependencies.\\[2pt]
\path{allFin_true}, \path{anyFin_true}, \path{eval_correct} & Checked Boolean
folds and finite-domain evaluator correctness. The reported dependencies
are \texttt{propext}, and additionally \texttt{Quot.sound} for the evaluator.\\[2pt]
Rational error and clock lemmas & Checked perturbation invariant, block
halving, amplitude absorption, and the full explicit
\path{solver_budget} schedule. The dependency audit reports
\texttt{propext}, \texttt{Classical.choice}, and \texttt{Quot.sound}.\\[2pt]
New main analytic results & Human proofs with explicit arithmetic targets;
no new main theorem in this archive has been delivered as a literal
\texttt{Prf PRAAx} proof value.\\
\bottomrule
\end{tabular}
\end{center}

The arithmetic axioms in the third row are inherited from the standard
\texttt{Rat} library: even its elementary power and order lemmas report
those dependencies. Thus the arithmetic host development is not advertised
as choice-free. Its proposed \(\ipra\) implementation uses the
denominator-cleared natural inequalities below. Host metatheoretic axioms
and the axioms available to an object derivation are different questions.

The checked project imports only \texttt{Std} and
\path{Init.Data.Rat.Lemmas}; it has no Mathlib or ProveIt import.
All project sources and their transitive project dependencies are included.
The external prerequisite is the pinned Lean distribution with its standard
library. \path{lean/PROVENANCE.json} records the compiler, distribution
checksum, and inspected repository sources. The command
\texttt{bash replay.sh} in the \texttt{lean} directory builds the library,
prints the dependency audit, and checks the examples. The compiler itself
and disposable build files are not included in the archive.

\subsection{Definitions at the semantic and arithmetic boundaries}

The checked logical syntax is completely explicit.
\texttt{Env X} means \texttt{Nat -> X}; \texttt{push} places a newly bound
variable at index zero and shifts the previous environment. The relation
\texttt{Related R} compares environments pointwise.
\texttt{PosFormula Atom} has constructors for an atom, truth, falsity,
conjunction, disjunction, and the two first-order quantifiers.
\texttt{Holds} recursively interprets those constructors.
\texttt{positive\_transfer} takes maps \(X\to Y\) and \(Y\to X\)
choosing related points. Existential transfer uses the first map;
universal transfer uses the second. Different margins can be incorporated
into the source and target atomic interpretations. Neither connective nor
quantifier adds a further margin. The Boolean evaluator uses
\texttt{allFin} and \texttt{anyFin}, proved correct by induction on the
finite domain size.

These definitions are ordinary Lean types. For an arithmetic instance, a
host function parameter is replaced by a specified PR term or a finite
input table with its proved contract. The additional data for the analytic
modules are as follows; all displayed predicates have been defined
mathematically in the preceding sections.

\begin{description}[style=nextline,leftmargin=1em]
\item[Rational and dense codes.]
A rational code is a signed integer numerator and a positive natural
denominator. Gaussian rational codes are pairs of these. A metric basis
has a decidable domain of natural codes, a fixed rational distance program,
and derived metric laws. For a pre-Hilbert basis use rational vector
operations and a rational inner product. The finite radius predicate is
\[
 \operatorname{Ball}(v,w,r)\;:\Longleftrightarrow\;
 r\geq0\ \land\ \langle v-w,v-w\rangle\leq r^2,
 \tag{F1}\label{eq:formal-ball}
\]
with validity of the two vector codes included in its domain guard.
It does not require an exact rational value of a generally irrational norm.

\item[Names, moduli, and observation contracts.]
For a fixed producer \(p(a,n)\), \(\Reg_p(a)\) is the formula from
Definition~\ref{def:regular}, with universal precision variables and
explicit metric tests. \(\Eq_{p,q}(a,b)\) is the stated finer-index equality
test. A dense evaluator stores its fixed evaluation term, monotone demand
term, and joint continuity derivations. Net data comprise a positive size,
an indexed list of valid centers, and the covering formula
\eqref{eq:top-net}. An evaluator for a function family includes exactly the
three rational contracts \eqref{eq:top-name}--\eqref{eq:top-space}.

\item[Polynomial and formula codes.]
A polynomial is a finite coefficient list; multiplication is convolution,
differentiation sends coefficient \(p_{j+1}\) to \((j+1)p_{j+1}\), and
integration from zero sends \(p_j\) to coefficient \(p_j/(j+1)\) at
index \(j+1\). A fixed multisorted signature records arities, input/output
sorts, symbol evaluators, and moduli. Terms and continuous formulas are
finite trees with the constructors in Section~\ref{sec:cm-model}; finite
environments respect the recorded sorts. Finite restriction replaces each
quantifier by a minimum or maximum scan, preserving its scope.
\end{description}

In the object implementation, tuple coding collects numerical parameters
into the arguments of fixed \texttt{PTm} terms. Inequality predicates are
quoted as \texttt{Fm PTm} formulas, with positive denominators cleared.
The intended output is \(\Gamma\vdash[\mathrm{PRAAx}]\varphi\).
The context \(\Gamma\) retains every input law that has not already been
proved. Merely filling a semantic record with assumptions does not
discharge those hypotheses.

\subsection{Two fully specified quantifier-free induction targets}

The new scalar estimates already have explicit arithmetic invariants.
Write the contraction factor as \(q=a/b\), initial majorant as \(R=r/s\),
per-step error as \(p/t\), and reserve as \(e/f\), with natural numerators,
positive denominators, and \(a\leq b\). The stationary-noise condition is
\[
                    pbf\leq(b-a)et.
\]
Define a PR numerator sequence
\[
 U_0=rt,\qquad U_{j+1}=aU_j+psb^{j+1}.
\]
Then \(U_j/(stb^j)\) is the exact majorant recurrence. The induction formula
is the quantifier-free inequality
\[
              fU_j\leq rtf a^j+est b^j.
 \tag{F2}\label{eq:formal-noise-invariant}
\]
Its base case is immediate. The next-step slack is identically
\[
 a\bigl[rtf a^j+estb^j-fU_j\bigr]
       +sb^j\bigl[(b-a)et-pbf\bigr].
\]
Both bracketed slacks are nonnegative. The repository's polynomial
inequality machinery can therefore prove the step with those two explicit
weights; \texttt{PRA.qfInd} yields the invariant. A trace whose error
certificates satisfy the local recurrence is compared with this majorant
by another bounded induction. No induction on a Cauchy predicate is used.

For the clock, put \(\delta=d/b\), \(0<d\leq b\). The corresponding
natural-number invariant is
\[
             (b+jd)(b-d)^j\leq b^{j+1}.
 \tag{F3}\label{eq:formal-clock-invariant}
\]
The exact step identity is
\[
 (b-d)\bigl(b+(j+1)d\bigr)+(j+1)d^2=b(b+jd).
\]
Multiply by \((b-d)^j\), discard the nonnegative term, and apply the
previous invariant. If desired, use a separate nonnegative parameter
\(a\) with \(a+d=b\) to eliminate truncated subtraction from the identity.
When \(hd\geq b\), specialization and positive cancellation give
\(2(b-d)^h\leq b^h\), the halving step. The elementary induction
\(C\leq2^C\) supplies the remaining precision offsets.

The checked declaration \texttt{solver\_budget} verifies the resulting
rational statement:
\[
\begin{gathered}
 0\leq\delta\leq1,\quad h\delta\geq1,\quad2\lambda\leq\delta C,\\
 x_0\leq R,\qquad
 x_{j+1}\leq(1-\delta)x_j+2\lambda2^{-(n+3+C)}
 \quad\Longrightarrow\quad
 x_{h(n+R+3)}\leq2^{-n-2}.
\end{gathered}
\]
Here \(h,C,R,n\) are natural numbers and \(x_j\) is a rational majorant.
The Hilbert application must establish its local radius contract
\eqref{eq:formal-ball}; this checked scalar theorem does not by itself
formalize an infinite-dimensional operator.

\subsection{Formalizing the main constructions in dependency order}

\paragraph{Completion and categorical operations.}
First quote the metric and joint evaluator laws, prove their substitution
instances, and apply the explicit demand in
Theorem~\ref{thm:cm-extension}. To derive preservation of \(\Reg\), introduce
arbitrary precision variables, instantiate the joint law, and combine its
rational inequalities. To derive preservation of \(\Eq\), use the finer
input indices and rational separation. These proofs do not apply
quantifier-free induction to the quantified formulas \(\Reg\) or \(\Eq\).
Product and change-of-basis maps use the same rules with maxima of demands.
Only after these obligations are proved should the fixed numerical family
be passed to \texttt{SetoidInterpretation.Presentation} and
\texttt{SetoidInterpretation.sound} \cite{repoSetoid}.

\paragraph{Function-space approximation and gluing.}
Define the rational positive part, the weights \(w_i\), partial sums
\(W_\ell\), and coefficient reconstruction. Prove the denominator bound
\(W_n\geq r/2\) from a bounded covering witness before invoking division.
For a fixed query \(x\) and evaluation precision \(p\), the central
sum invariant is
\[
 \left|\sum_{i<\ell}w_i(x)
            \bigl(h_i-E(a,x,p)\bigr)\right|
       \leq\bigl(2^{-k-1}+2^{-p}\bigr)W_\ell(x).
\]
The coefficient and support cases are decidable rational comparisons.
At the final sum, division yields \eqref{eq:top-object-target}.
Mixed-radix list coding proves the cardinality bound independently.
Gluing uses the same nonnegative-sum invariants and the already supplied
local bounds. None of these steps needs a decision about which arbitrary
profile classes are inhabited.

\paragraph{Continuous formulas and finite restrictions.}
Implement rational \(\min\), \(\max\), and \(1-u\), then finite extrema
with a first attaining index. The scan invariant says that the current
value is attained and is a lower, respectively upper, bound for the
processed prefix; its bounded universal part is Booleanized.
Prove that two finite lists related by the stated atomic-error relation
have extrema within that error. External structural induction on the
\emph{fixed} formula then constructs the arithmetic derivation of
Theorem~\ref{thm:cm-backforth}. Net comparison gives the coherent name in
Theorem~\ref{thm:cm-values}. The checked \texttt{positive\_transfer} is the
logical core of this argument; the numeric extrema lemma is additional
work, not a result already contained in that Boolean theorem.

\paragraph{The Hilbert solver and the Dirichlet application.}
Prove the dense inner-product identities and the finite Cauchy--Schwarz
inequality. Express approximate contraction as the rational rule
\[
 \operatorname{Ball}(v,w,r)\Longrightarrow
 \operatorname{Ball}\bigl(F_n(v),F_\ell(w),
                    qr+\eta_n+\eta_\ell\bigr).
 \tag{F4}\label{eq:formal-local-contraction}
\]
Derive this rule from the finite operator contracts; auxiliary approximation
errors are removed by the explicit rational separation lemma. Its
application to the ball, two-run, and adjacent-state invariants (H12)--(H14)
uses the same quantifier-free majorant calculations as
\eqref{eq:formal-noise-invariant}. The output name and its residual are
then proved coherent before the semantic fixed point is introduced.

For the PDE, define \(\operatorname{Int}(p)=\sum_jp_j/(j+1)\),
\(P_0\) by its two rational endpoint equations, and
\(\operatorname{Energy}(p,q)=\operatorname{Int}(p'q')\).
Polynomial antiderivatives implement \(G\) exactly. The approximating
operator has derivative
\(a_kv'-\operatorname{Int}(a_kv')\) and zero initial value.
Coefficient identities prove the endpoints, mean-subtraction identity,
and Green identity. Finite Riemann estimates with the polynomial modulus
prove the required nonnegative-integral inequalities. These derive the
strong approximation, coercivity, and norm contracts needed by
Theorem~\ref{thm:hilbert-solver}. Extension of the final bilinear identity
uses its explicit error bounds; no general Riesz or weak compactness
theorem is an omitted prerequisite.

\paragraph{Complex derivatives and finite analytic covers.}
Use finite Gaussian-rational coefficient lists and the derivative operation
defined above. Prove the finite differentiated-geometric inequality
of Lemma~\ref{lem:derivative-geometric}, including the \(L<j\) zero case,
by its finite coefficient-counting identity. This supplies the uniform
tail and coefficient-perturbation bounds for all \(j\leq J\).
The constants \(K_j,D_j\) and the positive radius gaps are rational;
Booleanized maxima over \(j\leq J\) allocate one common precision.
The clock calculations reuse the same block-halving arithmetic.
Finite binomial expansion proves the quadratic difference-quotient remainder
(A5); substitution of arbitrarily fine polynomial observations identifies
the formal derivatives with the complex derivatives. Finally encode the
coefficient grid by mixed radix and quote the explicit count (A8).
This proves the finite \(C^J\) cover, without selecting a convergent
subsequence or assuming that its centers belong to the input family.

\subsection{Repository interfaces and the final acceptance criterion}

The following existing imports provide the intended object backend; they
are not dependencies of the standalone checked project:
\begin{center}
\small
\begin{tabular}{p{0.42\textwidth}p{0.48\textwidth}}
\toprule
Within \path{IPRADerivations/} & Intended use\\
\midrule
\path{Foundations/ObservationComposition} &
\path{compose_obs}: combine finite queries and their maximum demand.\\
\path{Foundations/FiniteReplay} &
\path{replay}: preserve a completed trace when the input table is extended.\\
\path{Foundations/SetoidInterpretation} &
\path{sound}: transport a fixed presented first-order theory after its
domain, congruence, and axiom obligations are proved.\\
\path{Foundations/ContractionIterates} &
Existing exact half-contraction induction; extend it with the stationary
noise invariant rather than silently treating it as already present.\\
\path{Foundations/DiscreteGronwall} &
Existing denominator-cleared recurrence proofs. Its growth-factor theorem
is not itself the new contraction-reserve theorem.\\
\path{Analysis/ErrorBudgets} &
Finite sums, products, triangle estimates, and separated reciprocals.\\
\bottomrule
\end{tabular}
\end{center}

A completed integration should export explicit \texttt{Prf PRAAx}
declarations for the finite approximation, coherence, and residual targets,
retain their actual hypotheses, and audit their transitive proof
dependencies. Semantic adequacy should then identify each endpoint with
the represented theorem. A successful host Lean build is valuable evidence
for the arithmetic and structural arguments; it becomes evidence of
object-\(\ipra\) provability only when those additional proof values have
actually been constructed and checked.
